% !TEX root = ../thesis.tex

\chapter{Ciele záverečnej práce}\label{ch:goals}

Hlavným cieľom mojej práce je vytvoriť knižnicu pre programovací jazyk \emph{C}, ktorá bude reprezentovať svet \emph{Robota Karla}, a pomocou ktorej bude možné \emph{Robota Karla} aj ovládať. Moje čiastkové ciele, ktoré by som chcel dosiahnuť pri implementácii práce, som zhrnul do nasledujúcich bodov:

\begin{enumerate}
    \item Vytvoriť knižnicu v~jazyku \emph{C} reprezentujúcu svet \emph{Robota Karla}. Program pre \emph{Robota Karla} bude reprezentovaný štandardným programom jazyka \emph{C} a jednotlivé \emph{Karlove} príkazy ako aj senzory budú reprezentované pomocou funkcií.

    \item Vytvoriť jednoduché a prehľadné textové rozhranie pre zobrazenie aktuálnej situácie, v~ktorej sa \emph{Robot Karel} nachádza.

    \item Vytvoriť štúdijné materiály a ukážkové príklady pre použitie knižnice vo výučbe.

    \item Rozšíriť pôvodné schopností \emph{Robota Karla} o~nové schopnosti. Tie budú dostupné pomocou samostatných knižničných modulov.

    \item Otestovať knižnicu a materiály na čo najväčšej vzorke študentov, ideálne priamo na študentoch prvého ročníka. Výsledky získané testovaním budú použité vo výučbe úvodných programátorských kurzov poskytovaných \emph{Katedrou počítačov a informatiky}.
\end{enumerate}

